\section{KNN填补}

\begin{definition}[KNN 插值]
    KNN 插值是一种基于邻居的插值方法。对于每个缺失值，找出与其最相似的K个样本，然后用这K个样本的特征均值来填充缺失值。
    KNNImputer 通过欧几里德距离矩阵寻找最近邻。
\end{definition}

\begin{tips}
    \begin{itemize}
        \item 优点：能够利用数据的局部特性进行填充。
        \item 缺点：计算量较大，K值的选择对结果影响较大。
    \end{itemize}
\end{tips}

每个样本的缺失值，都是从其他训练集中找到最近邻的平均值估算的，如果两个都不缺失的特征都相近，那么两个样本就相近。

在找到距离待预测样本 $x$ 最近的 $k$ 个邻居（记为集合 $N_k(x)$）之后，我们需要根据这些邻居的信息来做出决策。

\begin{definition}[KNN 分类决策]
    对于分类问题，决策规则通常是“多数表决法”。即预测样本 $x$ 的类别 $\hat{y}$ 为其 $k$ 个最近邻居中出现次数最多的类别。
    $$
        \hat{y} = \underset{c \in \mathcal{Y}}{\arg\max} \sum_{(x_i, y_i) \in N_k(x)} I(y_i = c)
    $$
    其中：
    \begin{itemize}
        \item $\mathcal{Y}$ 是所有可能的类别标签的集合。
        \item $c$ 是其中一个类别。
        \item $I(\cdot)$ 是指示函数，当条件成立时为 1，否则为 0。
        \item $\arg\max$ 表示取能让后面表达式值最大的那个类别 $c$。
    \end{itemize}
\end{definition}

\begin{definition}[KNN 回归决策]
    对于回归问题，决策规则通常是“平均法”。即预测样本 $x$ 的值为其 $k$ 个最近邻居的真实值的平均数。
    $$
        \hat{y} = \frac{1}{k} \sum_{(x_i, y_i) \in N_k(x)} y_i
    $$
    有时也会使用“加权平均法”，即根据距离的远近为邻居分配不同的权重（例如，距离越近，权重越大），然后进行加权平均。若权重 $w_i = 1/d(x, x_i)$，则公式为：
    $$
        \hat{y} = \frac{\sum_{(x_i, y_i) \in N_k(x)} w_i y_i}{\sum_{(x_i, y_i) \in N_k(x)} w_i}
    $$
\end{definition}

\begin{lstlisting}[style=python]
from sklearn.impute import KNNImputer
import numpy as np

X = [[1, 2, np.nan], [3, 4, 3], [np.nan, 6, 5], [8, 8, 7]]
np.array(X)
Out[16]: 
array([[ 1.,  2., nan],
       [ 3.,  4.,  3.],
       [nan,  6.,  5.],
       [ 8.,  8.,  7.]])

imputer = KNNImputer(n_neighbors=2)
imputer.fit_transform(X)
Out[14]: 
array([[1. , 2. , 4. ],
       [3. , 4. , 3. ],
       [5.5, 6. , 5. ],
       [8. , 8. , 7. ]])
\end{lstlisting}

